% Computer Science A --- SQL language, modeling, DDL/DML, and transactions (\input from \texttt{main.tex}).
% SQL fragments: \verb|\lstinline[style=sqlinline]| (see \texttt{SQLEXAM} in \texttt{main.tex}).
% Free-response bank \texttt{csau3} at the front is migrated from \texttt{cs-a/03-practice.tex} (SQL, JDBC, keys, joins, advanced SQL, transactions/indexes).

\runningheader{Computer Science A}{SQL \& transactions}{Page \thepage\ of \numpages}

\begingradingrange{csau3}

\uplevel{\par\textbf{SQL and relational data}\par}

\question[4]
What is SQL and why would we use this language?
\begin{solutionorbox}[1.25in]
  Structured Query Language: declarative language for defining, querying, and manipulating relational data. Use it to persist and retrieve data in RDBMS products with a standard, portable syntax.
\end{solutionorbox}

\question[4]
What are the SQL sublanguages and their purpose?
\begin{solutionorbox}[1.45in]
  Common split: DDL (define schema), DML (insert/update/delete/select data), DCL (grant/revoke privileges), TCL (transactions: commit/rollback/savepoint). Purposes vary by vendor but these buckets organize commands.
\end{solutionorbox}

\question[4]
What is a RDBMS?
\begin{solutionorbox}[1.05in]
  Relational Database Management System: software that stores data as tables with keys and constraints, enforces relational rules, and processes SQL.
\end{solutionorbox}

\question[4]
Describe relational database tables.
\begin{solutionorbox}[1.25in]
  Two-dimensional relations: named columns (attributes) and rows (tuples/records); each row is unique (often via a primary key); tables link via foreign keys.
\end{solutionorbox}

\question[4]
What are constraints and can you describe a few constraints?
\begin{solutionorbox}[1.45in]
  Rules enforced by the DBMS: e.g.\ \texttt{PRIMARY KEY}, \texttt{FOREIGN KEY}, \texttt{UNIQUE}, \texttt{NOT NULL}, \texttt{CHECK}. They protect data quality and relationships.
\end{solutionorbox}

\question[4]
Why would I use the \texttt{WHERE} clause?
\begin{solutionorbox}[1.05in]
  To filter rows in \texttt{SELECT}, \texttt{UPDATE}, or \texttt{DELETE} so only rows matching a predicate are read or changed.
\end{solutionorbox}

\question[4]
What are some operators that can be used in SQL?
\begin{solutionorbox}[1.35in]
  Comparison (\texttt{=}, \texttt{<>}, \texttt{<}, \texttt{LIKE}, \texttt{IN}, \texttt{BETWEEN}), logical (\texttt{AND}, \texttt{OR}, \texttt{NOT}), arithmetic, string ops, and \texttt{IS NULL}---plus set operators in some contexts.
\end{solutionorbox}

\newpage

\uplevel{\par\textbf{JDBC and application data access}\par}

\question[4]
What is the JDBC API and the benefits of using it?
\begin{solutionorbox}[1.35in]
  Java Database Connectivity: standard API for connecting to databases, executing SQL, and reading results. Benefits: vendor-neutral code, connection pooling integration, tooling, and alignment with Java types.
\end{solutionorbox}

\question[4]
What is the DAO Design Pattern and why should we use it?
\begin{solutionorbox}[1.35in]
  Data Access Object: separate persistence logic (SQL/JDBC) from domain/services. Easier testing (mocks), swapping storage, and keeping SQL changes localized.
\end{solutionorbox}

\question[4]
What information would you need in order to successfully connect to a database?
\begin{solutionorbox}[1.35in]
  JDBC URL, driver class, credentials (user/password), optional pool settings, schema name, and SSL options as required by the deployment.
\end{solutionorbox}

\question[4]
What is the difference between a Simple and Prepared JDBC statement?
\begin{solutionorbox}[1.35in]
  \texttt{Statement} sends fixed SQL text. \texttt{PreparedStatement} precompiles a template with \texttt{?} placeholders, binds parameters safely, improves reuse, and helps prevent SQL injection.
\end{solutionorbox}

\question[4]
What is SQL Injection and how can we prevent it using the JDBC?
\begin{solutionorbox}[1.45in]
  Attacker embeds malicious SQL via untrusted input. Prevent with parameterized queries (\texttt{PreparedStatement}), validation, least-privilege DB users, and avoiding string concatenation for SQL.
\end{solutionorbox}

\newpage

\uplevel{\par\textbf{Keys, design, and relationships}\par}

\question[4]
What is a foreign key?
\begin{solutionorbox}[1.1in]
  A column (or set) in one table referencing a primary (or unique) key in another, encoding a relationship and enabling joins with referential checks.
\end{solutionorbox}

\question[4]
What is referential integrity?
\begin{solutionorbox}[1.15in]
  DB-enforced consistency: foreign key values must match existing parent keys (or be null), so orphans are blocked or cascades apply per policy.
\end{solutionorbox}

\question[4]
What is normalization?
\begin{solutionorbox}[1.35in]
  Organizing tables to reduce redundancy and anomalies (typically through normal forms), separating facts so each depends on the key, whole key, and nothing but the key.
\end{solutionorbox}

\question[4]
What is multiplicity?
\begin{solutionorbox}[1.2in]
  Cardinality of a relationship between entities (e.g.\ one-to-one, one-to-many, many-to-many), often modeled with keys and junction tables for $n$:$m$ cases.
\end{solutionorbox}

\newpage

\uplevel{\par\textbf{Joins, views, and set operations}\par}

\question[4]
Describe what a join is and explain the different types of joins we can create.
\begin{solutionorbox}[1.65in]
  Combine rows from two (or more) tables on a related column. Common kinds: \texttt{INNER} (matches only), \texttt{LEFT}/\texttt{RIGHT}/\texttt{FULL} outer (keep non-matching sides with nulls), \texttt{CROSS} (cartesian).
\end{solutionorbox}

\question[4]
What is the difference between a join and a set operation?
\begin{solutionorbox}[1.25in]
  Joins combine columns from related tables horizontally on predicates. Set ops (\texttt{UNION}, \texttt{INTERSECT}, \texttt{EXCEPT}) combine result sets of compatible queries vertically by rows.
\end{solutionorbox}

\question[4]
What is a view and why is it useful?
\begin{solutionorbox}[1.25in]
  A named, stored query (virtual table): simplifies complex SQL, can enforce column-level access, and stabilizes an API over underlying schema changes (implementation varies by DB).
\end{solutionorbox}

\newpage

\uplevel{\par\textbf{SQL: control, grouping, and programmability}\par}

\question[4]
What is DCL?
\begin{solutionorbox}[1in]
  Data Control Language: privilege commands such as \texttt{GRANT} and \texttt{REVOKE} to manage who may access which database objects.
\end{solutionorbox}

\question[4]
What is TCL?
\begin{solutionorbox}[1.05in]
  Transaction Control Language: \texttt{COMMIT}, \texttt{ROLLBACK}, \texttt{SAVEPOINT}---boundaries that make a set of changes atomic and durable per policy.
\end{solutionorbox}

\question[4]
What is the difference between an aggregate and a scalar function?
\begin{solutionorbox}[1.2in]
  Aggregates collapse many rows (\texttt{COUNT}, \texttt{SUM}, \texttt{AVG}) often with \texttt{GROUP BY}. Scalar functions compute per row (\texttt{UPPER}, \texttt{LENGTH}, \texttt{COALESCE}).
\end{solutionorbox}

\question[4]
What is the \texttt{GROUP BY} clause used for?
\begin{solutionorbox}[1.05in]
  Partition rows sharing grouping expressions so aggregates apply per group rather than over the whole table.
\end{solutionorbox}

\question[4]
What is the \texttt{ORDER BY} clause used for?
\begin{solutionorbox}[1in]
  Sort result rows by one or more columns/expressions, ascending or descending, optionally with collation-specific rules.
\end{solutionorbox}

\question[4]
What is a stored procedure in SQL?
\begin{solutionorbox}[1.2in]
  Named routine stored in the database (procedural SQL), invoked by apps; can encapsulate logic, reduce round trips, and centralize validation (trade-offs: portability, testing).
\end{solutionorbox}

\question[4]
What is a trigger in SQL?
\begin{solutionorbox}[1.2in]
  Procedure fired automatically on \texttt{INSERT}/\texttt{UPDATE}/\texttt{DELETE} (or DDL in some systems) to enforce rules, audit, or denormalized maintenance.
\end{solutionorbox}

\newpage

\uplevel{\par\textbf{Transactions, functions, and indexes}\par}

\question[4]
What is a transaction?
\begin{solutionorbox}[1.05in]
  A logical unit of work: a sequence of operations committed as one atomic bundle or rolled back on failure, isolating concurrent sessions per isolation level.
\end{solutionorbox}

\question[4]
What are the properties of a transaction?
\begin{solutionorbox}[1.25in]
  ACID: Atomicity (all-or-nothing), Consistency (valid state), Isolation (concurrent effects controlled), Durability (committed data survives crashes).
\end{solutionorbox}

\question[4]
What is a function in SQL?
\begin{solutionorbox}[1.15in]
  Named routine that typically returns a single value (scalar or table-valued) and can be used in expressions; may be built-in or user-defined depending on the DBMS.
\end{solutionorbox}

\question[4]
What is an index in SQL?
\begin{solutionorbox}[1.2in]
  Auxiliary structure (often B-tree) on column(s) to speed lookups and some joins/sorts at the cost of storage and write overhead; unique indexes also enforce uniqueness.
\end{solutionorbox}

\endgradingrange{csau3}

\newpage

\begingradingrange{csasql}

\uplevel{\par\textbf{SQL language basics}\par}

\question[4]
What is SQL?
\begin{choices}
  \choice A database for storing and maintaining state for data
  \choice A tool used solely for statistical analysis
  \choice A mentality for creating databases to be as least redundant as possible
  \CorrectChoice A language for interacting with a database
\end{choices}
\begin{solution}
  SQL (\textbf{Structured Query Language}) is the standard language for working with relational databases: querying, updating, and (with DDL) defining structure.
\end{solution}

\question[4]
What is the syntax to alter a table in SQL?
\begin{choices}
  \CorrectChoice \lstinline[style=sqlinline]!ALTER TABLE table_name ADD column_name data_type;!
  \choice \lstinline[style=sqlinline]!ALTER table_name ADD column_name data_type;!
  \choice \lstinline[style=sqlinline]!ALTER TABLE ADD column_name data_type;!
  \choice \lstinline[style=sqlinline]!ALTER TABLE table_name ADD column_name;!
\end{choices}
\begin{solution}
  \texttt{ALTER TABLE} names the table, then \texttt{ADD} introduces the new column and its type; dialects may add \texttt{IF NOT EXISTS}, constraints, or \texttt{AFTER} placement.
\end{solution}

\question[4]
An SQL \lstinline[style=sqlinline]!INSERT! statement inserts \blank[5em] record(s) into any single table in a relational database.
\begin{choices}
  \choice One
  \choice Two
  \CorrectChoice One or more
  \choice None of the above
\end{choices}
\begin{solution}
  \texttt{INSERT\ldots VALUES} can insert one row or several row tuples; \texttt{INSERT\ldots SELECT} can insert many rows from a query.
\end{solution}

\question[4]
Which of the following is \emph{not} considered part of a database's schema?
\begin{choices}
  \choice Tables
  \choice Columns
  \choice Table relationships
  \CorrectChoice Rows
\end{choices}
\begin{solution}
  The \textbf{schema} describes structure (tables, columns, keys, constraints, relationships). \textbf{Rows} are data instances, not structural metadata.
\end{solution}

\question[4]
Which of the following is true about DQL (Data Query Language)?
\begin{choices}
  \choice It is used to modify database structure
  \CorrectChoice It is used to retrieve data from a database
  \choice It controls access to database objects
  \choice It manages transactions within the database
\end{choices}
\begin{solution}
  DQL centers on \texttt{SELECT} (and related read-only retrieval). Structure changes are DDL; privilege changes are DCL; \texttt{COMMIT}/\texttt{ROLLBACK} are TCL.
\end{solution}

\question[4]
What does DDL stand for?
\begin{choices}
  \CorrectChoice Data Definition Language
  \choice Database Definition Language
  \choice Data Description Language
  \choice Data Define Language
\end{choices}
\begin{solution}
  \textbf{DDL} defines or changes schema: \texttt{CREATE}, \texttt{ALTER}, \texttt{DROP}, \texttt{TRUNCATE} (classification varies slightly by textbook/DBMS).
\end{solution}

\question[4]
What are the SQL sublanguages?
\begin{choices}
  \CorrectChoice DDL, DML, DQL, DCL, and TCL
  \choice DDL, DML, TQL, DCL, and RCL
  \choice TDL, DGL, DQL, DCL, and PCL
  \choice RDL, DML, DQL, QCL, and TCL
\end{choices}
\begin{solution}
  Common teaching split: \textbf{DDL} (schema), \textbf{DML} (row changes), \textbf{DQL} (\texttt{SELECT}), \textbf{DCL} (grants/revokes), \textbf{TCL} (transactions).
\end{solution}

\question[4]
Which join returns all records from the \emph{second} (right-hand) table, with matching rows from the other side when possible?
\begin{choices}
  \choice Inner join
  \choice Left join
  \CorrectChoice Right join
  \choice Self join
\end{choices}
\begin{solution}
  \texttt{RIGHT OUTER JOIN} keeps every row from the right operand; non-matches get \texttt{NULL}s on the left side. (Some teams prefer to rewrite as \texttt{LEFT JOIN} for consistency.)
\end{solution}

\question[4]
Which join returns all records from \emph{both} tables, filling with \texttt{NULL}s where there is no match?
\begin{choices}
  \choice Self join
  \choice Cross join
  \choice Natural join
  \CorrectChoice Full (outer) join
\end{choices}
\begin{solution}
  \texttt{FULL OUTER JOIN} is the usual answer; not every DBMS supports it (e.g.\ MySQL historically used workarounds).
\end{solution}

\question[4]
Which keyword introduces the join predicate between tables?
\begin{choices}
  \choice \lstinline[style=sqlinline]!IN!
  \CorrectChoice \lstinline[style=sqlinline]!ON!
  \choice \lstinline[style=sqlinline]!AT!
  \choice \lstinline[style=sqlinline]!AND!
\end{choices}
\begin{solution}
  \texttt{JOIN\ldots ON condition} specifies how rows relate; \texttt{AND} can extend the predicate but does not replace \texttt{ON}.
\end{solution}

\newpage

\question[4]
A foreign key must reference a primary key.
\begin{choices}
  \choice TRUE
  \CorrectChoice FALSE
\end{choices}
\begin{solution}
  In practice a foreign key usually references a \textbf{primary key} or another \textbf{UNIQUE} candidate key on the parent table---not strictly ``primary key only'' in the general SQL model.
\end{solution}

\question[4]
In an RDBMS, a foreign key is\ldots
\begin{choices}
  \choice A value used to uniquely identify a row
  \choice A virtual column which shows data from another table
  \CorrectChoice A constraint: a column (or columns) whose values reference a key in another table
  \choice Data loaded from an external database
\end{choices}
\begin{solution}
  A \textbf{foreign key} enforces referential integrity: child values must match an existing parent key (or be \texttt{NULL} if allowed).
\end{solution}

\question[4]
What is a foreign key?
\begin{choices}
  \choice The identifier for the database
  \choice A column or group of columns that uniquely identify a row in the same table
  \choice A key used to access a secured database table
  \CorrectChoice A column that references a column of another table to establish a relationship between rows
\end{choices}
\begin{solution}
  Foreign keys model parent/child relationships; the DBMS can enforce actions such as \texttt{ON DELETE CASCADE} or \texttt{SET NULL}.
\end{solution}

\question[4]
What is a primary key?
\begin{choices}
  \choice A constraint that requires uniqueness only
  \choice A constraint that requires a logical check only
  \CorrectChoice A constraint that combines \texttt{UNIQUE} and \texttt{NOT NULL}
  \choice A constraint that requires a reference to pre-existing data from another column
\end{choices}
\begin{solution}
  A table has at most one primary key; its columns are non-null and uniquely identify rows (often backed by a unique index).
\end{solution}

\question[4]
The \lstinline[style=sqlinline]!AS! keyword specifies an alias for a column or table.
\begin{choices}
  \CorrectChoice TRUE
  \choice FALSE
\end{choices}
\begin{solution}
  Example: \lstinline[style=sqlinline]!SELECT first_name AS fname FROM employees e!---aliases simplify reading and joins.
\end{solution}

\question[4]
An alias is an alternate name for a column or table to make it easier to reference.
\begin{choices}
  \CorrectChoice TRUE
  \choice FALSE
\end{choices}
\begin{solution}
  Aliases are especially useful in self-joins and when expressions appear in the \texttt{SELECT} list.
\end{solution}

\question[4]
What is the purpose of SQL constraints?
\begin{choices}
  \CorrectChoice To restrict which values may appear in a column (or across columns)
  \choice A limitation to the size of a table
  \choice To limit which operations a user may perform
  \choice A cap on how many operations the database may run per time period
\end{choices}
\begin{solution}
  Examples: \texttt{NOT NULL}, \texttt{UNIQUE}, \texttt{PRIMARY KEY}, \texttt{FOREIGN KEY}, \texttt{CHECK}, and \texttt{DEFAULT}.
\end{solution}

\question[4]
In PostgreSQL, \texttt{SERIAL} is\ldots
\begin{choices}
  \CorrectChoice A pseudotype that sets up an auto-incrementing integer column (via a sequence)
  \choice A constraint that column data must be serializable for transactions
  \choice A datatype meant only for human-readable serial numbers
  \choice An isolation level vulnerable to phantom reads
\end{choices}
\begin{solution}
  \texttt{SERIAL} is convenient syntax; modern PostgreSQL style often prefers identity columns (\texttt{GENERATED AS IDENTITY}) for clarity.
\end{solution}

\question[4]
Which constraint lets you declare a custom boolean condition on column values?
\begin{choices}
  \CorrectChoice \lstinline[style=sqlinline]!CHECK!
  \choice \texttt{SPECIAL}
  \choice \texttt{NOT NULL} (as the only answer)
  \choice \texttt{UNIQUE} (as the only answer)
\end{choices}
\begin{solution}
  \texttt{CHECK (condition)} rejects inserts/updates that violate the predicate, \\
  e.g.\ \lstinline[style=sqlinline]!CHECK (age >= 0)!.
\end{solution}

\newpage

\question[4]
Which is a valid \texttt{INSERT} statement?
\begin{choices}
  \choice \lstinline[style=sqlinline]|INSERT INTO employee(id,firstname,lastname) VALUES[75,'Edgar','Stephens'];|
  \choice \lstinline[style=sqlinline]!INSERT INTO TABLE employee(id,firstname,lastname) VALUES(48,'Adam','Parker');!
  \CorrectChoice \lstinline[style=sqlinline]!INSERT INTO employee(id,firstname,lastname) VALUES(39,'Margeret','Sampson');!
  \choice None of the above
\end{choices}
\begin{solution}
  Standard SQL uses parentheses for \texttt{VALUES}, not square brackets; \texttt{TABLE} does not belong between \texttt{INTO} and the table name.
\end{solution}

\question[4]
Which of the following is \emph{not} a conventional SQL sublanguage?
\begin{choices}
  \choice DML --- Data Manipulation Language
  \choice TCL --- Transaction Control Language
  \choice DQL --- Data Querying Language
  \CorrectChoice DPL --- Data Persistence Language
\end{choices}
\begin{solution}
  ``DPL'' is a distractor; persistence is an application/ORM concern, not a standard fifth SQL sublanguage label.
\end{solution}

\question[4]
DML statements include commands like \texttt{CREATE}, \texttt{DROP}, and \texttt{ALTER}.
\begin{choices}
  \choice TRUE
  \CorrectChoice FALSE
\end{choices}
\begin{solution}
  \texttt{CREATE}/\texttt{DROP}/\texttt{ALTER} are \textbf{DDL}. \textbf{DML} is typically \texttt{INSERT}, \texttt{UPDATE}, \texttt{DELETE}, and sometimes \texttt{SELECT} depending on the textbook.
\end{solution}

\question[4]
Which statement retrieves all columns and all rows from \texttt{Customers}?
\begin{choices}
  \CorrectChoice \lstinline[style=sqlinline]!SELECT * FROM Customers;!
  \choice \lstinline[style=sqlinline]!SELECT all FROM Customers;!
  \choice \lstinline[style=sqlinline]!SELECT % FROM Customers;!
  \choice \lstinline[style=sqlinline]!SELECT all columns FROM Customers;!
\end{choices}
\begin{solution}
  \texttt{*} projects every column; there is no standard \texttt{ALL} keyword for that purpose. \texttt{\%} is for \texttt{LIKE} patterns, not column lists.
\end{solution}

\question[4]
Which SQL sublanguage is used to create and change \emph{rows} in tables?
\begin{choices}
  \CorrectChoice DML
  \choice DQL
  \choice DCL
  \choice DDL
\end{choices}
\begin{solution}
  Row-level changes are \textbf{DML}; defining tables is \textbf{DDL}; reading with \texttt{SELECT} is often labeled \textbf{DQL}.
\end{solution}

\question[4]
What does SQL stand for?
\begin{choices}
  \choice Standalone Query Language
  \choice Standalone Query Lambdas
  \CorrectChoice Structured Query Language
  \choice Structured Query Lambdas
\end{choices}
\begin{solution}
  Originally ``SEQUEL'' historically; today \textbf{Structured Query Language} is the expansion students memorize.
\end{solution}

\question[4]
Which SQL sublanguage defines tables, schema, and databases?
\begin{choices}
  \choice DML
  \choice DQL
  \choice DCL
  \CorrectChoice DDL
\end{choices}
\begin{solution}
  \texttt{CREATE TABLE}, \texttt{CREATE DATABASE}, \texttt{ALTER}, and \texttt{DROP} are classic DDL examples.
\end{solution}

\question[4]
Which statement removes all rows from \texttt{table1} where \texttt{FirstName} is \texttt{'Mike'}?
\begin{choices}
  \choice \lstinline[style=sqlinline]!DROP FROM table1 WHERE FirstName = 'Mike';!
  \CorrectChoice \lstinline[style=sqlinline]!DELETE FROM table1 WHERE FirstName = 'Mike';!
  \choice \lstinline[style=sqlinline]!DROP FirstName = 'Mike' FROM table1;!
  \choice \lstinline[style=sqlinline]!DELETE FirstName = 'Mike' FROM table1;!
\end{choices}
\begin{solution}
  \texttt{DELETE FROM\ldots WHERE\ldots} removes matching rows. \texttt{DROP} drops objects (tables, views), not row predicates in this form.
\end{solution}

\question[4]
Which clause sorts the result of a \texttt{SELECT}?
\begin{choices}
  \choice \lstinline[style=sqlinline]!HAVING!
  \CorrectChoice \lstinline[style=sqlinline]!ORDER BY!
  \choice \lstinline[style=sqlinline]!GROUP BY!
  \choice \lstinline[style=sqlinline]!WHERE!
\end{choices}
\begin{solution}
  \texttt{ORDER BY} controls presentation order; 
  \texttt{WHERE} filters before grouping; \texttt{HAVING} filters grouped results.
\end{solution}

\question[4]
What does this query return?
\par\noindent\lstinline[style=sqlinline]!SELECT * FROM table1 WHERE firstname = 'Brad' OR firstname = 'Ana';!
\begin{choices}
  \choice All rows except the first \texttt{Brad} or \texttt{Ana} match
  \choice It is not valid SQL
  \choice Rows that are neither \texttt{Brad} nor \texttt{Ana}
  \CorrectChoice Only rows whose \texttt{firstname} is \texttt{Brad} or \texttt{Ana}
\end{choices}
\begin{solution}
  The \texttt{OR} combines two equality predicates; \texttt{*} returns all columns for matching rows.
\end{solution}

\newpage

\uplevel{\par\textbf{Grouping, keys, DDL, and modeling}\par}

\question[4]
Which clause is used with aggregate functions to partition rows into groups?
\begin{choices}
  \choice \texttt{ORDER BY}
  \CorrectChoice \texttt{GROUP BY}
  \choice \texttt{HAVING}
  \choice \texttt{WHERE}
\end{choices}
\begin{solution}
  \texttt{GROUP BY} collapses rows sharing the same grouping keys so aggregates (\texttt{COUNT}, \texttt{SUM}, \ldots) apply per group. \texttt{ORDER BY} only sorts; \texttt{HAVING} filters \emph{after} grouping; \texttt{WHERE} filters before grouping and cannot reference aggregates directly.
\end{solution}

\question[4]
Which two constraints are implicit when a column is declared \texttt{PRIMARY KEY}?
\begin{choices}
  \choice \texttt{UNIQUE} and \texttt{DEFAULT}
  \choice \texttt{NOT NULL} and \texttt{CHECK}
  \CorrectChoice \texttt{UNIQUE} and \texttt{NOT NULL}
  \choice \texttt{FOREIGN KEY} and \texttt{UNIQUE}
\end{choices}
\begin{solution}
  Primary keys must be non-null and unique; \texttt{DEFAULT} and arbitrary \texttt{CHECK} rules are not implied unless you add them.
\end{solution}

\question[4]
What is the correct syntax to add column \texttt{email} \texttt{VARCHAR(100)} to table \texttt{users}?
\begin{choices}
  \CorrectChoice \lstinline[style=sqlinline]!ALTER TABLE users ADD email VARCHAR(100);!
  \choice \lstinline[style=sqlinline]!MODIFY TABLE users ADD email VARCHAR(100);!
  \choice \lstinline[style=sqlinline]!UPDATE users ADD COLUMN email VARCHAR(100);!
  \choice \lstinline[style=sqlinline]!INSERT INTO users (email) VALUES (VARCHAR(100));!
\end{choices}
\begin{solution}
  \texttt{ALTER TABLE\ldots ADD} changes schema. \texttt{UPDATE} changes data; \texttt{INSERT} adds rows, not columns.
\end{solution}

\question[4]
What is the main difference between \texttt{TRUNCATE} and \texttt{DROP}?
\begin{choices}
  \choice \texttt{TRUNCATE} is DML and \texttt{DROP} is DDL
  \choice \texttt{DROP} removes only rows; \texttt{TRUNCATE} removes the table object
  \CorrectChoice \texttt{TRUNCATE} removes all rows but keeps the table; \texttt{DROP} removes the table object (and its rows)
  \choice They are synonyms
\end{choices}
\begin{solution}
  Both are destructive at different levels: \texttt{TRUNCATE} fast-empty (permissions and logging differ by engine); \texttt{DROP TABLE} removes definition and data. Classification as DDL vs DML varies slightly by product---focus on structural effect.
\end{solution}

\question[4]
What is the primary purpose of a \texttt{FOREIGN KEY} constraint?
\begin{choices}
  \choice To uniquely identify each row in one table
  \choice To forbid \texttt{NULL} in a column
  \choice To auto-increment integers
  \CorrectChoice To link rows in two tables and enforce referential integrity
\end{choices}
\begin{solution}
  The child column(s) must reference existing parent key values (subject to \texttt{ON DELETE}/\texttt{ON UPDATE} actions).
\end{solution}

\question[4]
When a parent row referenced by a child foreign key is deleted, what does \texttt{ON DELETE CASCADE} do?
\begin{choices}
  \choice Prevents deleting the parent
  \choice Sets the child's foreign key columns to \texttt{NULL}
  \CorrectChoice Deletes matching child rows automatically
  \choice Always raises a foreign-key error
\end{choices}
\begin{solution}
  \texttt{CASCADE} propagates the delete; \texttt{SET NULL} and \texttt{RESTRICT}/\texttt{NO ACTION} describe other common policies.
\end{solution}

\question[4]
Which command is a classic \textbf{DDL} example?
\begin{choices}
  \choice \texttt{SELECT}
  \choice \texttt{INSERT}
  \CorrectChoice \texttt{CREATE}
  \choice \texttt{UPDATE}
\end{choices}
\begin{solution}
  \texttt{CREATE} defines schema objects. \texttt{SELECT} is DQL; \texttt{INSERT}/\texttt{UPDATE} are DML in the usual split.
\end{solution}

\question[4]
In PostgreSQL, which type is a common shortcut for an auto-incrementing integer primary key?
\begin{choices}
  \choice \texttt{AUTO\_INCREMENT}
  \choice \texttt{IDENTITY}
  \choice \texttt{VARCHAR}
  \CorrectChoice \texttt{SERIAL}
\end{choices}
\begin{solution}
  \texttt{AUTO\_INCREMENT} is MySQL-flavored; \texttt{IDENTITY} is common in SQL Server; PostgreSQL traditionally uses \texttt{SERIAL} or identity columns.
\end{solution}

\question[4]
Sort primarily by \texttt{last\_name} ascending, using \texttt{first\_name} descending as a tie-breaker.
\begin{choices}
  \CorrectChoice \lstinline[style=sqlinline]!ORDER BY last_name ASC, first_name DESC!
  \choice \lstinline[style=sqlinline]!ORDER BY last_name, first_name!
  \choice \lstinline[style=sqlinline]!SORT BY last_name ASC, first_name DESC!
  \choice \lstinline[style=sqlinline]!GROUP BY last_name ASC, first_name DESC!
\end{choices}
\begin{solution}
  \texttt{ASC}/\texttt{DESC} apply per sort key; omitted direction defaults to ascending. \texttt{SORT BY} is not standard SQL; \texttt{GROUP BY} aggregates rather than sorts the final projection.
\end{solution}

\question[4]
In \texttt{SELECT}, which symbol selects all columns?
\begin{choices}
  \choice \texttt{\%}
  \CorrectChoice \texttt{*}
  \choice \texttt{\#}
  \choice \texttt{\&}
\end{choices}
\begin{solution}
  \texttt{*} means ``all columns'' in the \texttt{SELECT} list; \texttt{\%} is a wildcard for \texttt{LIKE} patterns.
\end{solution}

\question[4]
Which is a graphical representation of a database's entities and relationships?
\begin{choices}
  \choice Class and table structure only
  \choice Entity Framework
  \CorrectChoice Entity Relationship Diagram (ERD)
  \choice A generic configuration file
\end{choices}
\begin{solution}
  An \textbf{ERD} shows entities, attributes, and cardinalities; ORM tools may import from or export to ERDs.
\end{solution}

\question[4]
An ERD illustrates entities and the relationships between them.
\begin{choices}
  \CorrectChoice TRUE
  \choice FALSE
\end{choices}
\begin{solution}
  ERDs support logical/physical modeling before \texttt{CREATE TABLE} scripts.
\end{solution}

\question[4]
Views are normally used to make databases more complex to use.
\begin{choices}
  \CorrectChoice FALSE
  \choice TRUE
\end{choices}
\begin{solution}
  Views simplify access (hiding joins, enforcing column subsets) and can centralize security for read-only roles.
\end{solution}

\question[4]
Creating a view is done with which sublanguage?
\begin{choices}
  \CorrectChoice DDL
  \choice DML
  \choice DCL
  \choice TCL
\end{choices}
\begin{solution}
  \lstinline[style=sqlinline]!CREATE VIEW! defines a stored query object---schema/metadata, hence DDL in the usual curriculum.
\end{solution}

\newpage

\uplevel{\par\textbf{Database normalization}\par}

\question[4]
In the \blank[5em] normal form, it is ensured that a record's data for a specific column is atomic.
\begin{choices}
  \CorrectChoice First
  \choice Second
  \choice Third
  \choice Fourth
\end{choices}
\begin{solution}
  \textbf{First Normal Form (1NF)} requires atomic (indivisible) column values---no repeating groups or multi-valued cells in the relational model taught in most courses.
\end{solution}

\question[4]
What is a unique key?
\begin{choices}
  \choice A unique key is one or more columns that do \emph{not} uniquely identify a row in a table
  \CorrectChoice A unique key is one or more columns that uniquely identify a row in a table
  \choice Both are correct
  \choice None of the above
\end{choices}
\begin{solution}
  A \textbf{UNIQUE} constraint (or unique key) forbids duplicate combinations of values, like a primary key but you may allow multiple \texttt{NULL}s depending on the DBMS.
\end{solution}

\question[4]
What is 1NF?
\begin{choices}
  \CorrectChoice A table is referred to as being in First Normal Form if the column values are atomic
  \choice The table should not possess partial dependency
  \choice Normal Form ensures the reduction of data duplication
  \choice NONE
\end{choices}
\begin{solution}
  1NF is the ``atomic columns'' requirement; higher normal forms add further dependency rules on top of 1NF.
\end{solution}

\question[4]
A table is in \textbf{Second Normal Form (2NF)} when which of the following holds?
\begin{choices}
  \choice The table is in 1NF and no non-key attribute depends on another non-key attribute (no transitive dependencies)
  \CorrectChoice The table is in 1NF and every non-key attribute is fully functionally dependent on the whole primary key (no partial dependency on part of the key)
  \choice The table is in 1NF and column values are atomic only
  \choice The table removes all many-to-many relationships by using foreign keys
\end{choices}
\begin{solution}
  2NF rules out \textbf{partial} dependencies on a \textbf{composite} key: every non-key attribute must depend on the \emph{entire} key, not just a subset.
\end{solution}

\question[4]
In the context of 2NF, a \textbf{partial dependency} means:
\begin{choices}
  \choice Two different rows share the same foreign key value
  \CorrectChoice A non-key attribute can be determined from only \emph{part} of a composite primary key
  \choice A primary key is made of more than one column
  \choice A column allows \texttt{NULL} values
\end{choices}
\begin{solution}
  Partial dependency is the classic 2NF violation: some attribute depends on only a \textbf{proper subset} of the composite key.
\end{solution}

\question[4]
A table is in \textbf{Third Normal Form (3NF)} when which of the following holds?
\begin{choices}
  \choice The table is in 1NF and column values are atomic
  \choice The table is in 2NF and every non-key attribute depends only on candidate keys (Boyce\textendash Codd wording)
  \CorrectChoice The table is in 2NF and there are no transitive dependencies of non-key attributes on the primary key
  \choice The table is in 2NF and there are no multi-valued attributes left
\end{choices}
\begin{solution}
  3NF removes \textbf{transitive} chains where a non-key attribute determines another non-key attribute (unless the determinant is a superkey/candidate key, depending on definition).
\end{solution}

\question[4]
Which situation is a classic example of a \textbf{transitive dependency} (relevant to 3NF)?
\begin{choices}
  \choice A composite primary key has two columns
  \CorrectChoice A non-key attribute determines another non-key attribute (e.g.\ zip $\rightarrow$ city when neither is the primary key)
  \choice Two unrelated many-to-many relationships are stored in one table
  \choice A column stores multiple phone numbers separated by commas
\end{choices}
\begin{solution}
  If \texttt{zip} functionally determines \texttt{city} and both are non-keys, \texttt{city} is transitively dependent on the key through \texttt{zip}---a textbook 3NF violation.
\end{solution}

\question[4]
A table is in \textbf{Fourth Normal Form (4NF)} when which of the following holds (typical definition)?
\begin{choices}
  \choice It is in 3NF and every determinant is a candidate key
  \CorrectChoice It is in Boyce\textendash Codd Normal Form (BCNF) and independent multi-valued facts about the same entity are removed (no problematic multi-valued dependencies)
  \choice It is in 2NF and partial dependencies are removed
  \choice It is in 1NF and every attribute depends on the key, the whole key, and nothing but the key
\end{choices}
\begin{solution}
  4NF targets \textbf{multi-valued dependencies} (independent ``multi-valued facts'') beyond BCNF; wording varies by text but the MVD story is the core idea.
\end{solution}

\question[4]
Which issue is 4NF mainly designed to address that 3NF does \emph{not} fully cover?
\begin{choices}
  \choice Non-atomic column values (repeating groups)
  \choice Partial dependency on part of a composite key
  \choice Transitive dependency among non-key attributes
  \CorrectChoice Independent multi-valued relationships among attributes (multi-valued dependencies)
\end{choices}
\begin{solution}
  3NF/BCNF focus on functional dependencies; 4NF addresses cases where one key value is associated with \textbf{independent sets} of values that should be decomposed.
\end{solution}

\newpage

\uplevel{\par\textbf{Aggregates, transactions, and isolation}\par}

\question[4]
What is the difference between scalar functions and aggregate functions?
\begin{choices}
  \choice Aggregate functions return a single value, and scalar functions return a single value calculated from values in a column
  \CorrectChoice Scalar functions return a value per row; aggregate functions summarize many rows (often with \texttt{GROUP BY})
  \choice Aggregate functions cannot use a \texttt{GROUP BY} clause
\end{choices}
\begin{solution}
  \textbf{Scalar} functions (\texttt{UPPER}, \texttt{LENGTH}, \ldots) operate row-by-row. \textbf{Aggregates} (\texttt{COUNT}, \texttt{SUM}, \ldots) combine rows; they are commonly paired with \texttt{GROUP BY}.
\end{solution}

\question[4]
Which TCL command undoes work that has not been committed?
\begin{choices}
  \choice \texttt{CANCEL}
  \choice \texttt{REVERT}
  \CorrectChoice \texttt{ROLLBACK}
  \choice \texttt{ABORT}
\end{choices}
\begin{solution}
  \textbf{TCL} includes \texttt{COMMIT}, \texttt{ROLLBACK}, and often \texttt{SAVEPOINT}. \texttt{ROLLBACK} discards uncommitted changes (subject to engine rules).
\end{solution}

\question[4]
What is the use of the \texttt{SAVEPOINT} command from Transaction Control Language?
\begin{choices}
  \choice Load a file into the database
  \choice Update a file in the database
  \CorrectChoice Mark a point in a transaction to which you can roll back later
  \choice None of the above
\end{choices}
\begin{solution}
  \texttt{SAVEPOINT name} lets you \texttt{ROLLBACK TO SAVEPOINT name} without undoing the whole transaction.
\end{solution}

\question[4]
Larry loaded mismatched rows during an exam and time is running out. What is the quickest fix if nothing was committed?
\begin{choices}
  \choice \texttt{REVOKE}
  \choice \texttt{GRANT}
  \choice \texttt{ALTER}
  \CorrectChoice \texttt{ROLLBACK}
\end{choices}
\begin{solution}
  \texttt{ROLLBACK} drops uncommitted DML in the current transaction. Privilege and schema commands are the wrong tool for mistaken inserts.
\end{solution}

\newpage 


\question[4]
Several rows were inserted into the \texttt{Students} table in the current transaction. Which command makes those changes \textbf{durable} (persisted) so they survive connection loss?
\begin{choices}
  \choice \texttt{GRANT}
  \choice \texttt{SELECT}
  \CorrectChoice \texttt{COMMIT}
  \choice \texttt{ROLLBACK}
\end{choices}
\begin{solution}
  \texttt{COMMIT} ends the transaction and makes its effects permanent (until further DML). \texttt{ROLLBACK} discards uncommitted work. \texttt{GRANT}/\texttt{SELECT} do not end a DML transaction in this sense.
\end{solution}

\question[4]
What does the following do (assuming autocommit is off and a transaction is open)?
\begin{lstlisting}[style=sql]
DELETE FROM STUDENTS
WHERE AGE = 16;
ROLLBACK;
\end{lstlisting}
\begin{choices}
  \CorrectChoice Undoes the \texttt{DELETE} within the transaction (no durable removal if nothing was committed before \texttt{ROLLBACK})
  \choice Permanently deletes rows where \texttt{AGE = 16}
  \choice Drops the entire \texttt{STUDENTS} table
  \choice None of the above
\end{choices}
\begin{solution}
  \texttt{DELETE} removes matching rows in the transaction; \texttt{ROLLBACK} aborts the transaction and restores the pre-transaction state for uncommitted work (engine-dependent details aside, this is the textbook story).
\end{solution}

\question[4]
ACID stands for\ldots
\begin{choices}
  \CorrectChoice Atomicity, Consistency, Isolation, Durability
  \choice Authentic, Consistency, Immutable, Duplicable
  \choice Atomicity, Compatable, Immutable, Durability
  \choice Automated, Compatable, Isolation, Deterministic
\end{choices}
\begin{solution}
  \textbf{A}tomicity, \textbf{C}onsistency, \textbf{I}solation, \textbf{D}urability describe guarantees transactions aim to provide.
\end{solution}

\question[4]
What are transactions?
\begin{choices}
  \choice Only when you add information to a table
  \choice Only when you make changes to tables
  \CorrectChoice Commonly taught as bundling DML against table data with commit boundaries
  \choice Only when you create or delete tables
\end{choices}
\begin{solution}
  A \textbf{transaction} is a sequence of DML steps against table rows that ends with \texttt{COMMIT} (make durable) or \texttt{ROLLBACK} (discard). Example: table \texttt{orders(order\_id, total)}---you \texttt{INSERT} a new order line, \texttt{UPDATE} the customer's loyalty points, and \texttt{UPDATE} stock; either \textbf{all} of those changes become visible to others after \texttt{COMMIT}, or \textbf{none} of them do after \texttt{ROLLBACK}. DDL may or may not follow the same rules depending on the DBMS; coursework usually stresses DML + commit boundaries.
\end{solution}

\question[4]
What is a dirty read?
\begin{choices}
  \choice Reading data without a primary key
  \choice Reading data that is not normalized
  \choice Reading data you lack permission to read
  \CorrectChoice Reading another transaction's uncommitted data
\end{choices}
\begin{solution}
  \textbf{Dirty read} means Txn~B read a value that was \textbf{never committed}---only present in another session's in-flight transaction.

  \medskip
  \noindent\textbf{Committed starting state} (\texttt{accounts}):

  \begin{center}
  \footnotesize
  \begin{tabular}{|r|r|}
    \hline
    \textbf{id} & \textbf{balance} \\
    \hline
    1 & 100 \\
    \hline
  \end{tabular}
  \end{center}

  \noindent\textbf{Interleaving} (A writes \texttt{50} but does \textbf{not} \texttt{COMMIT}; B reads while A still holds the change):

  \begin{center}
  \footnotesize
  \renewcommand{\arraystretch}{1.12}
  \begin{tabular}{|c|p{0.40\textwidth}|p{0.40\textwidth}|}
    \hline
    \textbf{Step} & \textbf{Txn A} & \textbf{Txn B} \\
    \hline
    1 &
    \texttt{UPDATE accounts SET balance = 50 WHERE id = 1;}\newline
    (\textbf{no} \texttt{COMMIT} yet) &
    (idle) \\
    \hline
    2 &
    (still uncommitted) &
    \texttt{SELECT balance FROM accounts WHERE id = 1;}\newline
    Sees \textbf{50} $\leftarrow$ \textbf{dirty} (A may still roll back) \\
    \hline
    3 &
    \texttt{ROLLBACK;} $\Rightarrow$ row is \textbf{100} again in the database &
    B already believed \texttt{50} was real \\
    \hline
  \end{tabular}
  \renewcommand{\arraystretch}{1}
  \end{center}

  \begin{center}
  \footnotesize
  \begin{tikzpicture}[
    x=1.02cm, y=0.8cm,
    ax/.style={-Stealth, semithick},
    bar/.style={draw, rounded corners=2pt, minimum height=0.48cm, align=center, inner xsep=3pt, inner ysep=2pt}
  ]
    \node[anchor=east, font=\scriptsize\bfseries] at (-0.15, 1.35) {Txn A};
    \node[anchor=east, font=\scriptsize\bfseries] at (-0.15, 0.35) {Txn B};
    \node[bar, fill=red!10, minimum width=3.0cm] (aup) at (1.85, 1.35) {\scriptsize \texttt{UPDATE} $\to$ 50 (open)};
    \node[bar, fill=black!8, minimum width=1.35cm] (brd) at (1.85, 0.35) {\scriptsize \texttt{SELECT}};
    \node[bar, fill=gray!15, minimum width=1.25cm] (arb) at (4.55, 1.35) {\scriptsize \texttt{ROLLBACK}};
    \draw[ax] (0,-0.55) -- (6.2,-0.55) node[right, font=\scriptsize] {time $\to$};
    \foreach \x in {0,2,4,6} \draw (\x,-0.55) -- (\x,-0.62);
    \draw[dashed, black!45] (aup.south) -- (1.85,-0.55);
    \draw[dashed, black!45] (brd.south) -- (1.85,-0.55);
    \draw[dashed, black!45] (arb.south) -- (4.55,-0.55);
    \node[anchor=north, align=left, font=\scriptsize, text width=0.92\linewidth] at (3.1,-0.95) {B's \texttt{SELECT} lines up under A's open \texttt{UPDATE}: B observes \textbf{uncommitted} work. After \texttt{ROLLBACK}, the durable value is still \texttt{100}.};
  \end{tikzpicture}
  \end{center}
\end{solution}

\question[4]
Which isolation level allows dirty reads?
\begin{choices}
  \choice \texttt{SERIALIZABLE}
  \choice \texttt{REPEATABLE READ}
  \choice \texttt{READ COMMITTED}
  \CorrectChoice \texttt{READ UNCOMMITTED}
\end{choices}
\begin{solution}
  Only \texttt{READ UNCOMMITTED} is the rung that \textbf{permits dirty reads} in the textbook ladder: the engine may show you another transaction's \textbf{uncommitted} row versions (like \texttt{balance = 50} in the diagram above).

  \medskip
  \noindent\textbf{How this fits the ladder} (same story as the phantom section): each step up hides more anomalies. At \texttt{READ COMMITTED} and above, a normal read is expected to see only \textbf{committed} data, so the dirty \texttt{SELECT} during A's open transaction should not happen.

  \begin{center}
  \footnotesize
  \renewcommand{\arraystretch}{1.15}
  \begin{tabular}{|l|c|}
    \hline
    \textbf{Isolation level} & \textbf{Dirty reads allowed?} \\
    \hline
    \texttt{READ UNCOMMITTED}  & \textbf{Yes} \\
    \texttt{READ COMMITTED}    & No \\
    \texttt{REPEATABLE READ}   & No \\
    \texttt{SERIALIZABLE}      & No \\
    \hline
  \end{tabular}
  \renewcommand{\arraystretch}{1}
  \end{center}
\end{solution}

\question[4]
In transaction theory, what is a \textbf{non-repeatable read}?
\begin{choices}
  \choice Reading another transaction's \textbf{uncommitted} data (a dirty read)
  \choice Seeing \textbf{new matching rows} in a repeated query because of concurrent \texttt{INSERT}s (a phantom read)
  \CorrectChoice Reading the \textbf{same row} twice in one transaction and getting \textbf{different values} because another transaction \textbf{committed} an update \textbf{between} your two reads
  \choice An error that happens only when a table has no primary key
\end{choices}
\begin{solution}
  \textbf{Non-repeatable read} = your \textbf{second} read of the \textbf{same row} (same key / same \texttt{WHERE}) returns a \textbf{different} value than the first, and \textbf{both} values were already \textbf{committed} at the times you read them. Another session changed the row and \texttt{COMMIT}ted in between---no ``uncommitted peek'' required (that would be a \textbf{dirty read}).

  \medskip
  \noindent\textbf{One row, stable identity} (\texttt{products}):

  \begin{center}
  \footnotesize
  \begin{tabular}{|l|r|}
    \hline
    \textbf{sku} & \textbf{price} \\
    \hline
    ABC & 10 \\
    \hline
  \end{tabular}
  \end{center}

  \noindent\textbf{Timeline} (Txn~A never reads uncommitted data from B; B's \texttt{UPDATE} is \textbf{committed} before A's second read):

  \begin{center}
  \footnotesize
  \renewcommand{\arraystretch}{1.12}
  \begin{tabular}{|c|p{0.40\textwidth}|p{0.40\textwidth}|}
    \hline
    \textbf{Step} & \textbf{Txn A (you)} & \textbf{Txn B} \\
    \hline
    1 &
    \texttt{SELECT price FROM products WHERE sku = 'ABC';}\newline
    Result: \textbf{10} &
    (idle) \\
    \hline
    2 &
    (still in the same transaction) &
    \texttt{UPDATE products SET price = 12 WHERE sku = 'ABC';}\newline
    \texttt{COMMIT;} \\
    \hline
    3 &
    Same \texttt{SELECT} again:\newline
    Result: \textbf{12} $\leftarrow$ first and second read \textbf{disagree} &
    (done) \\
    \hline
  \end{tabular}
  \renewcommand{\arraystretch}{1}
  \end{center}

  \begin{center}
  \footnotesize
  \begin{tikzpicture}[
    x=1.02cm, y=0.8cm,
    ax/.style={-Stealth, semithick},
    bar/.style={draw, rounded corners=2pt, minimum height=0.48cm, align=center, inner xsep=3pt, inner ysep=2pt}
  ]
    \node[anchor=east, font=\scriptsize\bfseries] at (-0.15, 1.35) {Txn A};
    \node[anchor=east, font=\scriptsize\bfseries] at (-0.15, 0.35) {Txn B};
    \node[bar, fill=blue!8, minimum width=1.45cm] (a1) at (1.0, 1.35) {\scriptsize read \texttt{10}};
    \node[bar, fill=orange!12, minimum width=2.05cm] (bup) at (3.15, 0.35) {\scriptsize \texttt{UPDATE}+\texttt{COMMIT}};
    \node[bar, fill=blue!8, minimum width=1.45cm] (a2) at (5.5, 1.35) {\scriptsize read \texttt{12}};
    \draw[ax] (0,-0.55) -- (6.6,-0.55) node[right, font=\scriptsize] {time $\to$};
    \foreach \x in {0,2,4,6} \draw (\x,-0.55) -- (\x,-0.62);
    \draw[dashed, black!45] (a1.south) -- (1.0,-0.55);
    \draw[dashed, black!45] (bup.south) -- (3.15,-0.55);
    \draw[dashed, black!45] (a2.south) -- (5.5,-0.55);
    \node[anchor=north, align=left, font=\scriptsize, text width=0.92\linewidth] at (3.3,-0.95) {Contrast: \textbf{Dirty read} = see another txn's \textbf{uncommitted} value. \textbf{Phantom} = repeated query returns \textbf{different row count} because of \texttt{INSERT}/\texttt{DELETE}. Here the \textbf{row count stays 1}; only the column \texttt{price} changes.};
  \end{tikzpicture}
  \end{center}
\end{solution}

\newpage

\question[4]
Which isolation level allows non-repeatable reads (in the classic textbook ladder)?
\begin{choices}
  \choice \texttt{SERIALIZABLE}
  \choice \texttt{REPEATABLE READ}
  \CorrectChoice \texttt{READ COMMITTED}
  \choice \texttt{READ UNCOMMITTED}
\end{choices}
\begin{solution}
  In the \textbf{classic ladder}, \texttt{READ COMMITTED} lets each statement see the \textbf{latest committed} row. That is exactly the \texttt{products} / \texttt{10} then \texttt{12} story in the previous question: two reads in one transaction can observe two different committed versions, so reads are \textbf{not repeatable}.
 
  \medskip
  \noindent\textbf{Where non-repeatable reads stop} (textbook model): at \texttt{REPEATABLE READ} and \texttt{SERIALIZABLE}, your transaction is supposed to see a \textbf{stable} snapshot of rows it already read---so the second \texttt{SELECT} would still show \texttt{10} until you end your transaction (details vary by product).

  \begin{center}
  \footnotesize
  \renewcommand{\arraystretch}{1.15}
  \begin{tabular}{|l|c|c|}
    \hline
    \textbf{Isolation level} & \textbf{Dirty read?} & \textbf{Non-repeatable read?} \\
    \hline
    \texttt{READ UNCOMMITTED}  & Yes & Yes \\
    \texttt{READ COMMITTED}    & No  & \textbf{Yes} \\
    \texttt{REPEATABLE READ}   & No  & No (row snapshot) \\
    \texttt{SERIALIZABLE}      & No  & No \\
    \hline
  \end{tabular}
  \renewcommand{\arraystretch}{1}
  \end{center}

  \noindent\footnotesize\textit{Exam tip:} phantoms (new rows) are a separate column in the full matrix; \texttt{REPEATABLE READ} can still allow phantoms in the pure ladder even when \textbf{row} values look stable.
\end{solution}


\newpage 

\question[4]
In transaction theory, what is a \textbf{phantom read}?
\begin{choices}
  \CorrectChoice Seeing \textbf{new rows} that match your query predicate because another transaction inserted (and committed) them while yours runs
  \choice Seeing new \textbf{tables} created during your transaction
  \choice Seeing \textbf{updated} values in rows you already read (non-repeatable read)
  \choice Seeing updated \textbf{table names} during your transaction
\end{choices}
\begin{solution}
  A \textbf{phantom} is about the \textbf{shape of a query result} (how many rows match), not ``this one cell changed'' on a row you already pinned down (that is closer to a \textbf{non-repeatable read}).

  \medskip
  \noindent\textbf{Snapshot of committed rows} (\texttt{major = 'CS'}) at the start:

  \begin{center}
  \footnotesize
  \begin{tabular}{|r|l|}
    \hline
    \textbf{id} & \textbf{major} \\
    \hline
    101 & CS \\
    102 & CS \\
    103 & CS \\
    \hline
  \end{tabular}
  \end{center}

  \noindent\textbf{Timeline} (Txn~A keeps running; Txn~B commits an \texttt{INSERT} in the middle):

  \begin{center} 
  \footnotesize
  \renewcommand{\arraystretch}{1.15}
  \begin{tabular}{|c|p{0.42\textwidth}|p{0.42\textwidth}|}
    \hline
    \textbf{Step} & \textbf{Txn A (you)} & \textbf{Txn B (other session)} \\
    \hline
    1 &
    \texttt{SELECT COUNT(*) FROM students WHERE major = 'CS';}\newline
    Result: \textbf{3} (three rows above) &
    (idle) \\
    \hline
    2 &
    (still open) &
    \texttt{INSERT INTO students (id, major) VALUES (104, 'CS');}\newline
    \texttt{COMMIT;} $\Rightarrow$ fourth committed CS row exists \\
    \hline
    3 &
    Same query again:\newline
    \texttt{SELECT COUNT(*) \ldots}\newline
    Result: \textbf{4} --- a \textbf{new matching row} ``appeared'' in your result &
    (done) \\
    \hline
  \end{tabular}
  \renewcommand{\arraystretch}{1}
  \end{center}

  \begin{center}
  \footnotesize
  \begin{tikzpicture}[
    x=1.05cm, y=0.85cm,
    every node/.style={font=\footnotesize},
    ax/.style={-Stealth, semithick},
    bar/.style={draw, rounded corners=2pt, minimum height=0.5cm, align=center, inner xsep=3pt, inner ysep=2pt}
  ]
    % Swimlanes: Txn A (top) and Txn B (middle); time axis at bottom (no stacked bars on one row).
    \node[anchor=east, font=\scriptsize\bfseries] at (-0.15, 1.35) {Txn A};
    \node[anchor=east, font=\scriptsize\bfseries] at (-0.15, 0.35) {Txn B};
    \node[bar, fill=black!6, minimum width=2.0cm] (a1) at (1.1, 1.35) {\scriptsize 1st \texttt{COUNT}$=3$};
    \node[bar, fill=black!6, minimum width=2.1cm] (a2) at (6.5, 1.35) {\scriptsize 2nd \texttt{COUNT}$=4$};
    \node[bar, fill=orange!12, minimum width=2.2cm] (b1) at (3.6, 0.35) {\scriptsize \texttt{INSERT}+\texttt{COMMIT}};
    \draw[ax] (0,-0.55) -- (8.2,-0.55) node[right, font=\scriptsize] {time $\to$};
    \foreach \x in {0,2,4,6,8} \draw (\x,-0.55) -- (\x,-0.62);
    \draw[dashed, black!45] (a1.south) -- (1.1,-0.55);
    \draw[dashed, black!45] (b1.south) -- (3.6,-0.55);
    \draw[dashed, black!45] (a2.south) -- (6.5,-0.55);
    \node[align=left, anchor=north, font=\scriptsize, text width=0.92\linewidth] at (4.1,-0.95) {The higher second count is a \textbf{phantom}: a new committed row matched your \texttt{WHERE} clause between the two counts.};
  \end{tikzpicture}
  \end{center}
\end{solution}


\newpage 

\question[4]
Which isolation level is associated with phantom reads in the classic SQL ladder?
\begin{choices}
  \choice \texttt{SERIALIZABLE}
  \CorrectChoice \texttt{REPEATABLE READ}
  \choice Snapshot Isolation
  \choice Read Stability
\end{choices}
\begin{solution}
  Courseware often draws isolation as a \textbf{ladder} of strength (weakest at the bottom):

  \begin{center}
  \footnotesize
  \begin{tikzpicture}[
    box/.style={draw, minimum width=3.6cm, minimum height=0.62cm, align=center},
    arr/.style={-Stealth, thick}
  ]
    \node[box, fill=black!5] (ru) at (0,0) {\texttt{READ UNCOMMITTED}};
    \node[box, fill=black!5] (rc) at (0,-1.0) {\texttt{READ COMMITTED}};
    \node[box, fill=blue!10] (rr) at (0,-2.0) {\texttt{REPEATABLE READ}};
    \node[box, fill=black!5] (se) at (0,-3.0) {\texttt{SERIALIZABLE}};
    \draw[arr] (ru.south) -- (rc.north);
    \draw[arr] (rc.south) -- (rr.north);
    \draw[arr] (rr.south) -- (se.north);
    \node[anchor=west, align=left] at (4.0, -1.5) {\textbf{Stronger}$\;\to\;$fewer anomalies\\\small (more locking or stricter checks)};
  \end{tikzpicture}
  \end{center}

  \noindent\textbf{Classic textbook anomaly matrix} (what \emph{may} still happen at each rung). Use this for the exam story; real products deviate (especially MySQL/InnoDB \texttt{REPEATABLE READ}):

  \begin{center}
  \footnotesize
  \renewcommand{\arraystretch}{1.2}
  \begin{tabular}{|l|c|c|c|}
    \hline
    \textbf{Isolation level} & \textbf{Dirty read?} & \textbf{Non-repeatable read?} & \textbf{Phantom read?} \\
    \hline
    \texttt{READ UNCOMMITTED}  & Yes & Yes & Yes \\
    \texttt{READ COMMITTED}    & No  & Yes & Yes \\
    \texttt{REPEATABLE READ}   & No  & No  & \textbf{Yes} \\
    \texttt{SERIALIZABLE}      & No  & No  & No \\
    \hline
  \end{tabular}
  \renewcommand{\arraystretch}{1}
  \end{center}

  \noindent So \textbf{phantom reads} are still ``on the table'' at \texttt{REPEATABLE READ} in the \textbf{pure ladder}; only \texttt{SERIALIZABLE} removes them in that simplified model. The earlier \texttt{students} / \texttt{COUNT(*)} picture is exactly a phantom: repeated scans can see \textbf{new committed rows} that match your predicate.
\end{solution}


\newpage 


\question[4]
Which isolation level is typically the \textbf{strictest} (strongest guarantees) and therefore often the \textbf{slowest} due to extra locking or equivalent mechanisms?
\begin{choices}
  \CorrectChoice \texttt{SERIALIZABLE}
  \choice \texttt{REPEATABLE READ}
  \choice \texttt{READ COMMITTED}
  \choice \texttt{READ UNCOMMITTED}
\end{choices}
\begin{solution}
  \texttt{SERIALIZABLE} aims for results as if transactions ran \textbf{one after another} (no unsafe interleavings). Concretely, consider one row: \texttt{tickets(event\_id, seats\_left)} with \texttt{seats\_left = 1}.

  \medskip
  \noindent\textbf{Table before any sale}

  \begin{center}
  \footnotesize
  \begin{tabular}{|l|r|}
    \hline
    \textbf{event\_id} & \textbf{seats\_left} \\
    \hline
    Concert-A & 1 \\
    \hline
  \end{tabular}
  \end{center}

  \noindent\textbf{Dangerous interleaving} under weaker isolation (both clerks ``see'' one seat and both sell):

  \begin{center}
  \footnotesize
  \renewcommand{\arraystretch}{1.15}
  \begin{tabular}{|c|p{0.42\textwidth}|p{0.42\textwidth}|}
    \hline
    \textbf{Step} & \textbf{Txn T1 (clerk 1)} & \textbf{Txn T2 (clerk 2)} \\
    \hline
    1 & \texttt{SELECT seats\_left} $\Rightarrow$ \textbf{1} & \\
    2 & & \texttt{SELECT seats\_left} $\Rightarrow$ \textbf{1} \\
    3 & \texttt{UPDATE tickets SET seats\_left = 0 \ldots};\ \texttt{COMMIT}; & \\
    4 & & \texttt{UPDATE tickets SET seats\_left = 0 \ldots};\ \texttt{COMMIT}; \\
    \hline
  \end{tabular}
  \renewcommand{\arraystretch}{1}
  \end{center}

  \noindent After step~4 the row shows \texttt{0}, but \textbf{two sales} were allowed---\textbf{oversell}. The diagram below shows the overlap of the two reads on the same committed value:

  \begin{center}
  \footnotesize
  \begin{tikzpicture}[
    x=1.0cm, y=0.75cm,
    ax/.style={-Stealth, thick},
    every node/.style={font=\footnotesize}
  ]
    % T1 and T2 on separate tracks; bars are sequential in time so rectangles do not overlap.
    \node[anchor=east, font=\scriptsize\bfseries] at (-0.2, 1.45) {T1};
    \node[anchor=east, font=\scriptsize\bfseries] at (-0.2, 0.35) {T2};
    \draw[fill=blue!12, draw=black, rounded corners=1pt] (0.0,1.25) rectangle (1.9,1.65) node[midway, font=\scriptsize] {read $=1$};
    \draw[fill=blue!12, draw=black, rounded corners=1pt] (4.0,1.25) rectangle (5.4,1.65) node[midway, font=\scriptsize] {write $0$};
    \draw[fill=orange!15, draw=black, rounded corners=1pt] (1.5,0.15) rectangle (3.4,0.55) node[midway, font=\scriptsize] {read $=1$};
    \draw[fill=orange!15, draw=black, rounded corners=1pt] (5.2,0.15) rectangle (6.8,0.55) node[midway, font=\scriptsize] {write $0$};
    \draw[ax] (0,-0.55) -- (7.4,-0.55) node[right, font=\scriptsize] {time};
    \foreach \x/\lbl in {0/1,1.5/2,4/3,5.2/4} \draw (\x,-0.55) -- (\x,-0.62) node[below, font=\tiny, yshift=-1pt] {\lbl};
    \draw[dashed, black!40] (0.95,1.25) -- (0.95,-0.55);
    \draw[dashed, black!40] (2.45,0.15) -- (2.45,-0.55);
    \draw[dashed, black!40] (4.7,1.25) -- (4.7,-0.55);
    \draw[dashed, black!40] (6.0,0.15) -- (6.0,-0.55);
    \node[anchor=north, font=\scriptsize, text width=0.92\linewidth, align=left] at (3.7,-1.15) {Dashed drops mark the ordering from the step table: both reads see \texttt{1}, then both commits write \texttt{0}---\textbf{oversell}.};
  \end{tikzpicture}
  \end{center}

  \noindent\footnotesize\textbf{SERIALIZABLE} blocks unsafe interleavings like the one above: the engine enforces a serializable schedule (waits, conflict detection, retries or aborts). That extra synchronization is why it is often the \textbf{slowest} rung in practice.
\end{solution}


\newpage 

\uplevel{\par\textbf{Tables, SQL classification, JDBC, and design patterns}\par}

\question[4]
Which of the following terms is used to describe records in a database?
\begin{choices}
  \CorrectChoice Rows
  \choice Columns
  \choice Attributes
  \choice Properties
\end{choices}
\begin{solution}
  A \textbf{row} (tuple) is one record; \textbf{columns} (often called attributes) are fields of that record.
\end{solution}

\question[4] 
\textit{(Select all that apply.)} A column declared \lstinline[style=sqlinline]!PRIMARY KEY! is implicitly\ldots
\begin{checkboxes}
  \CorrectChoice \lstinline[style=sqlinline]!UNIQUE!
  \choice \lstinline[style=sqlinline]!FINAL!
  \choice \lstinline[style=sqlinline]!SERIAL! 
  \choice \lstinline[style=sqlinline]!FOREIGN KEY!
  \CorrectChoice \lstinline[style=sqlinline]!NOT NULL!
\end{checkboxes}
\begin{solution}
  A primary key must identify rows uniquely and cannot be null; hence it implies \textbf{UNIQUE} and \textbf{NOT NULL}. \texttt{FINAL}, \texttt{SERIAL}, and \texttt{FOREIGN KEY} are unrelated distractors (exact enforcement details vary slightly by DBMS).
\end{solution}


\newpage


\question[4]
\textit{(Select all that apply.)} Which of the following assertions about \texttt{CHAR}, \texttt{VARCHAR}, and \texttt{TEXT} are \emph{incorrect}?
\begin{checkboxes}
  \choice Strings inserted into a \texttt{CHAR} column will be padded, while strings inserted into a \texttt{VARCHAR} column will not be.
  \choice The \texttt{TEXT} type has a limited size according to the SQL standard, but can hold a string of any arbitrary size in PostgreSQL.
  \CorrectChoice A length can be defined for \texttt{CHAR}, \texttt{VARCHAR}, and \texttt{TEXT} columns.
  \CorrectChoice A length must be defined for \texttt{CHAR} and \texttt{VARCHAR} columns.
\end{checkboxes}
\begin{solution}
  The first two bullets are \textbf{correct} facts (so do \emph{not} select them). You should mark the last two bullets: they over-generalize how lengths work across \texttt{CHAR}, \texttt{VARCHAR}, and \texttt{TEXT}.

  \medskip\noindent\textbf{Portable DDL shape (PostgreSQL-style).} Fixed-width \texttt{CHAR(\textit{n})} and capped \texttt{VARCHAR(\textit{n})} declare a length in parentheses; \texttt{TEXT} is written \emph{without} \texttt{(\textit{n})}.
\begin{lstlisting}[style=sql]
CREATE TABLE labels (
  country_code CHAR(2),        -- always two logical characters; may pad on store/compare
  city         VARCHAR(64),   -- up to 64 characters; no CHAR-style padding
  blurb        TEXT            -- unbounded string type in PG; no VARCHAR-style "(n)" here
);
\end{lstlisting}
  In PostgreSQL, \texttt{TEXT(\textit{n})} is not valid syntax, so ``you can always slap \texttt{(\textit{n})} on \texttt{TEXT} like \texttt{VARCHAR}'' is a bad mental model.

  \medskip\noindent\textbf{Padding vs variable width (conceptual ``tape'' view).} Same logical insert \texttt{'AB'}:
\begin{center}
{\ttfamily\small
\begin{tabular}{@{}r@{ }l@{}}
  \texttt{CHAR(5)}     & \texttt{[A][B][\_][\_][\_]} \quad $\leftarrow$ fixed width 5 (pads) \\
  \texttt{VARCHAR(5)}  & \texttt{[A][B]} \quad $\leftarrow$ variable length up to cap 5 (no pads) \\
\end{tabular}\par}
\end{center}
\begin{lstlisting}[style=sql]
-- same idea as literal strings (illustrative):
-- CHAR(5)    -> 'AB   '
-- VARCHAR(5) -> 'AB'
\end{lstlisting}

  \medskip\noindent\textbf{Why ``\texttt{VARCHAR} \emph{must} declare a length'' is false (product-dependent).} PostgreSQL happily accepts \texttt{VARCHAR} with no \texttt{(\textit{n})} (unlimited length, same storage class as \texttt{TEXT} in practice):
\begin{lstlisting}[style=sql]
CREATE TABLE ex (
  title VARCHAR   -- legal in PostgreSQL: no "(n)" required
);
\end{lstlisting}
  Other engines (for example SQL Server) often require \texttt{VARCHAR(\textit{n})} or \texttt{VARCHAR(MAX)} instead, so there is no single universal ``must declare length'' rule.

  \medskip\noindent\textbf{On the \texttt{TEXT} bullet you left unselected.} SQL's large character types and vendor \texttt{TEXT} types have implementation limits even when the type name suggests ``big string''; PostgreSQL still cannot store an infinitely long value in one row. The exam contrast is about \textbf{type syntax and portability}, not infinite storage.
\end{solution}

\newpage 


\question[4]
What is the purpose of a \texttt{ResultSet}?
\begin{choices}
  \choice To create new records
  \CorrectChoice To use data retrieved from the database
  \choice To convert a row from a table into a Java \texttt{Set}
  \choice To check if there are duplicates in a table
\end{choices}
\begin{solution}
  A \texttt{ResultSet} is the JDBC cursor for reading query results (\texttt{next()}, getters, etc.). Inserts/updates use \texttt{executeUpdate} (or similar), not the result set itself.
\end{solution}

\question[4]
Which JDBC interface is used to create an instance of a \texttt{Statement}?
\begin{choices}
  \CorrectChoice \texttt{Connection}
  \choice \texttt{Statement}
  \choice \texttt{PreparedStatement}
  \choice \texttt{ResultSet}
\end{choices}
\begin{solution}
  Call \texttt{conn.createStatement()} or \texttt{conn.prepareStatement(\ldots)} on a \texttt{Connection} (from \texttt{DriverManager} or a \texttt{DataSource}). The UML-style map below matches the official JDBC API structure (core \texttt{java.sql}/\texttt{javax.sql} types):

  \begin{center}
  \begin{tikzpicture}[
      interface/.style={
          rectangle,
          draw,
          minimum width=3.0cm, 
          minimum height=0.95cm,
          align=center,
          font=\ttfamily\small
      },
      class/.style={
          rectangle,
          draw,
          line width=1.0pt,
          minimum width=3.0cm,
          minimum height=0.95cm,
          font=\ttfamily\small
      },
      extends/.style={
          ->,
          >=Triangle,
          thick
      },
      interface-extends/.style={
          ->,
          >=Triangle,
          thick,
          dashed
      },
      uses/.style={
          ->
      }
  ]
    % Row 1: connection entry points
    \node[class] (drivermanager) at (-0.2,0) {DriverManager};
    \node[interface] (connection) at (6.6,0) {$\ll$interface$\gg$\\\textit{Connection}};

    % Row 2: statement family root + result output
    \node[interface] (statement) at (2.5,-2.4) {$\ll$interface$\gg$\\\textit{Statement}};
    \node[interface] (resultset) at (-3.0,-2.4) {$\ll$interface$\gg$\\\textit{ResultSet}};

    % Row 3: specialized statements
    \node[interface] (prepared) at (2.5,-4.8) {$\ll$interface$\gg$\\\textit{PreparedStatement}};
    \node[interface] (callable) at (2.5,-7.2) {$\ll$interface$\gg$\\\textit{CallableStatement}};

    % Factory / usage relations (dashed like temp2's interface links)
    \draw[uses] (drivermanager.east) -- node[above, font=\scriptsize, align=center] {$\ll$create$\gg$\\\texttt{getConnection()}} (connection.west);
    \draw[uses] (connection.south east) to[out=250, in=20] node[pos=0.28, above left, font=\scriptsize, align=center] {$\ll$create$\gg$\\\texttt{createStatement()}} (statement.east);
    \draw[uses] (connection.south east) to[out=265, in=20] node[left, font=\scriptsize, align=center] {$\ll$create$\gg$\\\texttt{prepareStatement()}} (prepared.east);
    \draw[uses] (connection.south east) to[out=285, in=20] node[left, font=\scriptsize, align=center] {$\ll$create$\gg$\\\texttt{prepareCall()}} (callable.east);
    \draw[uses] (statement) -- node[above, font=\scriptsize, align=center] {$\ll$create$\gg$\\\texttt{executeQuery()}} (resultset);

    % Interface inheritance (PreparedStatement extends Statement, CallableStatement extends PreparedStatement)
    \draw[interface-extends] (prepared) -- (statement);
    \draw[interface-extends] (callable) -- (prepared);
  \end{tikzpicture}
  \end{center}
\end{solution}

\question[4]
\texttt{DriverManager} is the first interface you need to use in order to set up JDBC functionality.
\begin{choices}
  \choice TRUE
  \CorrectChoice FALSE
\end{choices}
\begin{solution}
  \texttt{DriverManager} is a \textbf{class}, not an interface; JDBC also does not mandate it as the literal first step (e.g.\ \texttt{DataSource} is common in applications).
\end{solution}

\question[4]
In SQL, which \lstinline[style=sqlinline]!JOIN! will return the matching rows plus the ones that are null in the first table?
\begin{choices}
  \choice \lstinline[style=sqlinline]!INNER JOIN!
  \CorrectChoice \lstinline[style=sqlinline]!LEFT JOIN!
  \choice \lstinline[style=sqlinline]!RIGHT JOIN!
  \choice \lstinline[style=sqlinline]!CROSS JOIN!
\end{choices}
\begin{solution}
  \texttt{LEFT OUTER JOIN} keeps every row from the left (first) table; non-matches get \texttt{NULL}s from the right side. (The stem follows the original quiz wording: ``null in the first table'' means nulls \emph{filled in} for missing right-hand columns.)
\end{solution}



\question[4]
What is the difference between \lstinline[style=sqlinline]!INNER JOIN! and \lstinline[style=sqlinline]!OUTER JOIN!?
\begin{choices}
  \choice An outer join will return only the rows that match based on the join predicate.
  \choice An outer join will return all rows in the left table and those that match in the right table.
  \choice An inner join will return rows that are not common between both tables.
  \CorrectChoice Inner join will return only the rows that match based on the join predicate.
\end{choices}
\begin{solution}
  An \textbf{inner join} keeps only rows satisfying the \texttt{ON} predicate. \textbf{Outer} joins also preserve non-matching rows from the left, right, or both sides (with \texttt{NULL}s), depending on \texttt{LEFT}/\texttt{RIGHT}/\texttt{FULL}.
\end{solution}

\question[4]
In the statement \lstinline[style=sqlinline]!SELECT * FROM users!, what's the \lstinline[style=sqlinline]!*! referring to?
\begin{choices}
  \CorrectChoice Columns
  \choice Rows
  \choice Tables
  \choice Records
\end{choices}
\begin{solution}
  \texttt{*} projects all \textbf{columns}; it does not mean ``all rows''---row filtering is \texttt{WHERE}/\texttt{HAVING}.
\end{solution}

\question[4]
The \lstinline[style=sqlinline]!UPDATE! operation is categorized under which SQL sublanguage?
\begin{choices}
  \CorrectChoice DML
  \choice DQL
  \choice DDL
  \choice DCL
\end{choices}
\begin{solution}
  \texttt{UPDATE} changes stored rows---classic \textbf{DML} alongside \texttt{INSERT} and \texttt{DELETE}. (Some courses label \texttt{SELECT} as \textbf{DQL}.)
\end{solution}

\newpage

\question[4]
\textit{(Select all that apply.)} Select the operations which are a part of DDL.
\begin{checkboxes}
  \choice \lstinline[style=sqlinline]!SELECT!
  \CorrectChoice \lstinline[style=sqlinline]!CREATE!
  \CorrectChoice \lstinline[style=sqlinline]!ALTER!
  \choice \lstinline[style=sqlinline]!UPDATE!
  \CorrectChoice \lstinline[style=sqlinline]!TRUNCATE!
\end{checkboxes}
\begin{solution}
  \texttt{CREATE}, \texttt{ALTER}, and \texttt{TRUNCATE} are commonly grouped with \textbf{DDL} in coursework (along with \texttt{DROP}). \texttt{SELECT} and \texttt{UPDATE} are not DDL. Classification of \texttt{TRUNCATE} varies by text/DBMS but matches the source quiz here.
\end{solution}
 
\newpage

\question[4]
What is the difference between the simple and prepared statement?
\begin{lstlisting}[style=java]
Connection conn = getConnection();
String userInput = getInput(); 
Statement st = conn.createStatement();
st.executeQuery("SELECT * FROM " + userInput);
PreparedStatement ps = conn.prepareStatement("SELECT * FROM ?");
ps.setString(1, userInput);
ps.executeQuery();
\end{lstlisting}
\begin{choices}
  \CorrectChoice The \texttt{PreparedStatement} prevents SQL injection by parameterizing and pre-compiling the query string.
  \choice There is no difference
  \choice The \texttt{Statement} cannot take user input, but \texttt{PreparedStatement} can.
  \choice Only \texttt{PreparedStatement} is a part of JDBC; \texttt{Statement} is not.
\end{choices}
\begin{solution}
  The key idea in the course material is \textbf{parameterization} vs string concatenation for untrusted input (\textbf{SQL injection}). The sample \texttt{SELECT * FROM ?} is not valid portable JDBC for binding a \emph{table} name---placeholders are for values; use a fixed SQL shape or validated identifiers separately.
\end{solution}

\question[4]
Which of the following are used in the Data Access Object design pattern?
\begin{choices}
  \CorrectChoice Abstraction via interfaces
  \choice Singleton pattern via private constructor
  \choice Polymorphism via method overloading
  \choice Recursion
\end{choices}
\begin{solution}
  DAOs centralize persistence behind an API (often an \textbf{interface}) so callers stay decoupled from JDBC/SQL details. The other options are unrelated distractors.
\end{solution}

\question[4]
Identify the correct syntax to retrieve all attributes of the \texttt{employees} table who have a salary over \$50{,}000 and do not work in \texttt{Sales}:
\begin{choices}
  \choice \lstinline[style=sqlinline]!QUERY id FROM employees GIVEN salary > 50000 AND dept != 'Sales';!
  \CorrectChoice \lstinline[style=sqlinline]!SELECT * FROM employees WHERE salary > 50000 AND dept <> 'Sales';!
  \choice \lstinline[style=sqlinline]!SELECT * FROM employees HAVING salary GREATER THAN 50000 OR dept <> 'Sales';!
  \choice \lstinline[style=sqlinline]!QUERY FROM employees WHERE salary > 50000 AND dept != 'Sales';!
\end{choices}
\begin{solution}
  Row filters use \texttt{WHERE}; \texttt{HAVING} applies after \texttt{GROUP BY}. Use standard \\
  \texttt{SELECT\ldots FROM\ldots WHERE}. (SQL string literals use single quotes; \texttt{<>} or \texttt{!=} for ``not equal'' depends on dialect---both are common.)
\end{solution}

\newpage


\question[4]
Which of the following is \emph{not} a command in the DML sublanguage of SQL?
\begin{choices}
  \choice \lstinline[style=sqlinline]!INSERT INTO!
  \choice \lstinline[style=sqlinline]!UPDATE!
  \choice \lstinline[style=sqlinline]!DELETE!
  \CorrectChoice \lstinline[style=sqlinline]!QUERY!
\end{choices}
\begin{solution}
  \texttt{INSERT}, \texttt{UPDATE}, and \texttt{DELETE} are classic \textbf{DML}. Retrieval is \texttt{SELECT}; \texttt{QUERY} is not a SQL keyword in that sense.
\end{solution}

\question[4]
Which of the following is \emph{not} a command in the DDL sublanguage of SQL?
\begin{choices}
  \choice \lstinline[style=sqlinline]!CREATE!
  \CorrectChoice \lstinline[style=sqlinline]!DELETE!
  \choice \lstinline[style=sqlinline]!ALTER!
  \choice \lstinline[style=sqlinline]!DROP!
\end{choices}
\begin{solution}
  \texttt{DELETE} removes rows (\textbf{DML}); \texttt{CREATE}, \texttt{ALTER}, and \texttt{DROP} are classic \textbf{DDL}.
\end{solution}

\newpage

\uplevel{\par\textbf{JDBC (quiz bank)}\par}

\question[4]
In JDBC, results from an SQL query are returned in\ldots
\begin{choices}
  \CorrectChoice \texttt{ResultSet}
  \choice \texttt{HashSet}
  \choice \texttt{SortedSet}
  \choice None of the above
\end{choices}
\begin{solution}
  \texttt{executeQuery} returns a \texttt{ResultSet} cursor over rows; collection types are not the JDBC query return type.
\end{solution}

\question[4]
In JDBC, in \texttt{PreparedStatement}, data is replaced with\ldots
\begin{choices}
  \CorrectChoice Question marks (\texttt{?})
  \choice Variable names
  \choice Dollar signs
  \choice Exclamation points
\end{choices}
\begin{solution}
  Placeholders \texttt{?} are bound with \texttt{setString}, \texttt{setInt}, etc., keeping literals out of the SQL string.
\end{solution}

\question[4]
Which of these is \emph{not} an interface in JDBC?
\begin{choices}
  \choice \texttt{Connection}
  \choice \texttt{Statement}
  \choice \texttt{PreparedStatement}
  \choice \texttt{CallableStatement}
  \CorrectChoice \texttt{ExecutableStatement}
\end{choices}
\begin{solution}
  \texttt{java.sql} defines \texttt{Connection}, \texttt{Statement}, \texttt{PreparedStatement}, and \texttt{CallableStatement}; \texttt{ExecutableStatement} is a distractor name.
\end{solution}

\question[4]
Which of the following is correct about the \texttt{ResultSet} class of JDBC?
\begin{choices}
  \choice \texttt{ResultSet} holds data retrieved from a database after you execute an SQL query using \texttt{Statement} objects
  \choice It acts as an iterator to allow you to move through its data
  \choice The \texttt{java.sql.ResultSet} interface represents the result of a database query
  \CorrectChoice All of the above
\end{choices}
\begin{solution}
  \texttt{ResultSet} is an \textbf{interface} (the stem says ``class'' in some banks) for tabular query results; you traverse rows with \texttt{next()} and read columns with getters.
\end{solution}

\question[4]
In JDBC, which statement type is the usual choice to bind parameters with \texttt{?} and reduce SQL injection risk?
\begin{choices}
  \choice \texttt{Statement}
  \CorrectChoice \texttt{PreparedStatement}
  \choice \texttt{CallableStatement}
  \choice None of the above
\end{choices}
\begin{solution}
  \texttt{PreparedStatement} separates SQL structure from user-supplied data; never concatenate untrusted input into raw SQL strings.
\end{solution}

\question[4]
A \texttt{CallableStatement} protects a statement from SQL injection.
\begin{choices}
  \choice TRUE
  \CorrectChoice FALSE
\end{choices}
\begin{solution}
  \texttt{CallableStatement} is for calling procedures; safety still depends on how parameters are bound and how dynamic SQL is built. \texttt{PreparedStatement} is the usual teaching answer for bound parameters on arbitrary SQL.
\end{solution}

\question[4]
What does JDBC provide?
\begin{choices}
  \choice A way to automatically back up databases
  \choice A way to map tables to an ER diagram
  \choice A way to turn database tables into Plain Old Java Objects (POJOs)
  \CorrectChoice An API to communicate with a database
\end{choices}
\begin{solution}
  JDBC is the standard Java \textbf{API} for relational DB access; ORMs and backups are separate concerns.
\end{solution}

\question[4]
JDBC operations typically throw which type of exception?
\begin{choices}
  \choice \texttt{IOException}
  \choice \texttt{JDBCConnectionException}
  \CorrectChoice \texttt{SQLException}
  \choice \texttt{SQLSyntaxException}
\end{choices}
\begin{solution}
  Checked \texttt{SQLException} (and often subclasses) wraps database errors; \texttt{SQLSyntaxException} is a subclass in some drivers, not the umbrella type named in every textbook.
\end{solution}

\question[4]
Which of the following correctly executes a \textbf{query} (DQL) using JDBC?
\begin{choices}
  \CorrectChoice Call \texttt{executeQuery()} on a \texttt{Statement} (or \texttt{PreparedStatement})
  \choice Call \texttt{executeQuery()} on a \texttt{ResultSet}
  \choice Call \texttt{executeStatement()} on a \texttt{Statement}
  \choice Call \texttt{executeStatement()} on a \texttt{ResultSet}
\end{choices}
\begin{solution}
  \texttt{Statement}/\texttt{PreparedStatement} expose \texttt{executeQuery} for \texttt{SELECT}-style commands; there is no standard \texttt{executeStatement} on \texttt{ResultSet}.
\end{solution}

\question[4]
Which Java package contains the core JDBC interfaces and classes?
\begin{choices}
  \choice \texttt{java.lang}
  \CorrectChoice \texttt{java.sql}
  \choice \texttt{java.util}
  \choice \texttt{java.io}
\end{choices}
\begin{solution}
  Core JDBC types live in \texttt{java.sql}; \texttt{javax.sql} adds pooling, \texttt{DataSource}, and related APIs.
\end{solution}

\question[4]
What is the purpose of the \texttt{DriverManager} class?
\begin{choices}
  \choice To manage \texttt{Statement} objects
  \CorrectChoice To obtain connections to a database (via a JDBC URL and registered drivers)
  \choice To manage \texttt{ResultSet} objects
  \choice To run the database server process on your machine
\end{choices}
\begin{solution}
  \texttt{DriverManager.getConnection(url, user, password)} is the classic entry point in coursework; applications often prefer a \texttt{DataSource} instead.
\end{solution}

\question[4]
What does JDBC stand for?
\begin{choices}
  \choice Java Database Collector
  \choice Java Data Bean Connectivity
  \CorrectChoice Java Database Connectivity
  \choice Java Database Communicator
\end{choices}
\begin{solution}
  \textbf{J}ava \textbf{D}atabase \textbf{C}onnectivity: the bridge between Java and relational databases through drivers.
\end{solution}

\question[4]
You should use \texttt{PreparedStatement} over \texttt{Statement} when binding user input because \texttt{PreparedStatement} helps protect against SQL injection.
\begin{choices}
  \CorrectChoice TRUE
  \choice FALSE
\end{choices}
\begin{solution}
  Parameter binding avoids interpreting user text as SQL syntax; combine with validation and least privilege for defense in depth.
\end{solution}

\question[4]
You need a \texttt{ResultSet} for typical DML statements (\texttt{INSERT}, \texttt{UPDATE}, \texttt{DELETE}).
\begin{choices}
  \choice TRUE
  \CorrectChoice FALSE
\end{choices}
\begin{solution}
  JDBC uses two common execution paths: \textbf{read} vs \textbf{write}. Queries (\texttt{SELECT}) return rows via \texttt{executeQuery()} $\rightarrow$ \texttt{ResultSet}; DML typically uses \texttt{executeUpdate()} and returns an affected-row \texttt{int}.

  \begin{center}
  \footnotesize
  \renewcommand{\arraystretch}{1.15}
  \begin{tabular}{@{}lll@{}}
    \hline
    \textbf{SQL kind} & \textbf{Call} & \textbf{You get} \\
    \hline
    DQL (\texttt{SELECT}) & \texttt{executeQuery()} & \texttt{ResultSet} (iterate rows) \\
    DML (\texttt{INSERT}, \texttt{UPDATE}, \texttt{DELETE}) & \texttt{executeUpdate()} & \texttt{int} (affected row count) \\
    \hline
  \end{tabular}
  \renewcommand{\arraystretch}{1}
  \end{center}

  \noindent\textit{Example:} \texttt{executeUpdate("DELETE FROM flights WHERE status = 'cancelled'")} might return \texttt{5} if five rows were deleted---a row count, not the deleted rows as a \texttt{ResultSet}.

\begin{lstlisting}[style=java]
// DQL: ask for data -> ResultSet
ResultSet rs = ps.executeQuery(
    "SELECT id FROM flights WHERE status = 'open'");

// DML: change data -> int row count
int removed = ps.executeUpdate(
    "DELETE FROM flights WHERE status = 'cancelled'");
\end{lstlisting}

\begin{lstlisting}[style=java]
// Advanced: generated keys after INSERT (still executeUpdate first)
PreparedStatement ps = conn.prepareStatement(
    "INSERT INTO flights (code) VALUES (?)",
    java.sql.Statement.RETURN_GENERATED_KEYS);
ps.setString(1, "AA123");
int rows = ps.executeUpdate(); // still executeUpdate, not executeQuery
ResultSet newIds = ps.getGeneratedKeys(); // optional key(s)
\end{lstlisting}
\end{solution}

\question[4]
What layer should JDBC persistence logic usually live in?
\begin{choices}
  \choice Controller layer
  \CorrectChoice DAO (data access) layer
  \choice Service layer only
  \choice All of the above
\end{choices}
\begin{solution}
  The \textbf{DAO} isolates SQL/JDBC from UI and business rules; services orchestrate use cases and may call DAOs.
\end{solution}

\newpage

\uplevel{\par\textbf{JDBC and stored procedures}\par}

\question[4]
Stored procedures can be invoked from Java using \texttt{java.sql.CallableStatement}.
\begin{choices}
  \CorrectChoice TRUE
  \choice FALSE
\end{choices}
\begin{solution}
  \texttt{CallableStatement} is the JDBC interface for calling database procedures/functions (often with \texttt{\{call ...\}} style escape syntax). \texttt{PreparedStatement} is for parameterized SQL; \texttt{Statement} is the basic executor.
\end{solution}

\endgradingrange{csasql}
